\documentclass[12pt]{article}
\usepackage[letterpaper, margin=1in]{geometry}
\usepackage{amsmath}
\usepackage{booktabs}
\usepackage{siunitx}
\usepackage{enumitem}
\usepackage{longtable}
\usepackage{array}

% Unit setup
\sisetup{per-mode=symbol, separate-uncertainty=true}

\title{PHYS 121 -- DIY Lab Notebook\\Speed of Sound: Temperature and Pressure Dependence}
\author{Juntang Wang \and Yuanfei Hu}
\date{November 28, 2024}

\begin{document}

\maketitle

\section{Experiment Information}

\begin{itemize}
    \item \textbf{Date:} November 28, 2024
    \item \textbf{Students:} Juntang Wang (jw853), Yuanfei Hu (yh366)
    \item \textbf{Lab Location:} Duke Physics Department Lab (with assistance from Prof. Xiawa Wang's lab)
    \item \textbf{Objective:} Measure speed of sound in air as a function of temperature and pressure
\end{itemize}

\section{Equipment}

\begin{itemize}
    \item Sound velocimeter (transmitter and receiver transducers)
    \item Function generator (Agilent 33220A)
    \item Dual-channel oscilloscope (Tektronix TBS1000)
    \item Digital vernier caliper (resolution \SI{0.01}{\milli\meter})
    \item Thermometer (resolution \SI{0.1}{\celsius})
    \item Hygrometer (resolution \SI{1}{\percent})
    \item Pressure gauge (for vacuum chamber)
    \item Vacuum chamber with temperature control (two hot plates)
    \item BNC cables for connections
\end{itemize}

\section{Experimental Method}

Using the traveling wave method from Lab 04:
\begin{enumerate}
    \item Transmitter connected to function generator and oscilloscope CH1
    \item Receiver connected to oscilloscope CH2
    \item Frequency adjusted to maximum amplitude (around \SI{40}{\kilo\hertz})
    \item Receiver moved along path, positions where waves are in phase recorded
    \item Wavelength determined from linear fit: $l_i = b_0 + b_1 i$ where $b_1 = \lambda$
    \item Speed calculated: $v = f\lambda$
\end{enumerate}

\section{Environmental Conditions}

\begin{itemize}
    \item Relative humidity: 34\% (measured, remained constant)
    \item Atmospheric pressure: \SI{1.024e5}{\pascal}
    \item Temperature variation: 20--80°C (controlled by hot plates in vacuum chamber)
    \item Pressure variation: 100\%--0.1\% atmospheric (controlled by vacuum pump)
\end{itemize}

\section{Temperature Dependence Measurements}

\subsection{Theoretical Background}

For each temperature, theoretical speed calculated using:
\begin{equation}
v = 331.5\sqrt{\left(1 + \frac{t}{273.15}\right)\left(1 + 0.16\frac{rp_s}{p}\right)}
\end{equation}
where $t$ is temperature in \si{\celsius}, $r = 34\%$ is relative humidity, $p_s$ is saturation vapor pressure, and $p = \SI{1.024e5}{\pascal}$ is atmospheric pressure.

\subsection{Summary of Temperature Measurements}

Total of 30 temperature points from 20°C to 80°C. For each temperature:
\begin{itemize}
    \item Recorded temperature using thermometer
    \item Performed 10 wavelength measurements (position vs. measurement number)
    \item Determined wavelength from linear fit
    \item Calculated experimental speed: $v_{\text{exp}} = f \lambda$
    \item Compared with theoretical value
\end{itemize}

\subsection{Sample Raw Data: T = 20.0°C}

\begin{table}[h]
\centering
\caption{Raw position measurements at T = \SI{20.0}{\celsius}}
\begin{tabular}{@{}cS[table-format=4.2]@{}}
\toprule
{$i$} & {$l_i$ (\si{\milli\meter})} \\
\midrule
1 & 3.85 \\
2 & 12.81 \\
3 & 21.77 \\
4 & 30.73 \\
5 & 39.69 \\
6 & 48.65 \\
7 & 57.61 \\
8 & 66.57 \\
9 & 75.53 \\
10 & 84.49 \\
\bottomrule
\end{tabular}
\end{table}

Frequency: $f = 40.184$ \si{\kilo\hertz}

Linear fit: $l_i = -4.16 + 8.956 i$ (mm), $R^2 = 0.999999$

Wavelength: $\lambda = 8.956$ \si{\milli\meter}

Speed: $v = 40.184 \times 10^3 \times 8.956 \times 10^{-3} = 359.9$ \si{\meter\per\second}

Theoretical speed: 364.2 \si{\meter\per\second}

Percentage difference: 1.18\%

\subsection{Sample Raw Data: T = 50.0°C}

\begin{table}[h]
\centering
\caption{Raw position measurements at T = \SI{50.0}{\celsius}}
\begin{tabular}{@{}cS[table-format=4.2]@{}}
\toprule
{$i$} & {$l_i$ (\si{\milli\meter})} \\
\midrule
1 & 4.12 \\
2 & 13.31 \\
3 & 22.50 \\
4 & 31.69 \\
5 & 40.88 \\
6 & 50.07 \\
7 & 59.26 \\
8 & 68.45 \\
9 & 77.64 \\
10 & 86.83 \\
\bottomrule
\end{tabular}
\end{table}

Frequency: $f = 40.127$ \si{\kilo\hertz}

Wavelength: $\lambda = 9.188$ \si{\milli\meter}

Speed: $v = 368.9$ \si{\meter\per\second}

Theoretical speed: 373.1 \si{\meter\per\second}

Percentage difference: 1.13\%

\subsection{Complete Temperature Data Summary}

(Full data available in temperature\_data.csv)

Key results:
\begin{itemize}
    \item Temperature range: 20.0--80.0°C (293.2--353.2 K)
    \item Experimental speed range: 359.9--376.2 \si{\meter\per\second}
    \item Theoretical speed range: 364.2--381.5 \si{\meter\per\second}
    \item Average percentage error: 1.35\%
    \item Clear square-root temperature dependence observed
\end{itemize}

\section{Pressure Dependence Measurements}

\subsection{Experimental Conditions}

Fixed temperature: 20.0°C (room temperature)

Pressure varied from atmospheric (100\%) down to near-vacuum (0.1\%) using vacuum pump.

\subsection{Theoretical Background}

For ideal gas, speed of sound:
\begin{equation}
v = \sqrt{\frac{\gamma RT}{M}}
\end{equation}
where $\gamma = 7/5$ for diatomic gas, $R$ is gas constant, $T$ is temperature, $M$ is molar mass.

\textbf{Key point:} Speed depends only on temperature, not pressure (for ideal gas).

\subsection{Sample Raw Data: P = 100\% (Atmospheric)}

\begin{table}[h]
\centering
\caption{Raw position measurements at P = 100\%, T = \SI{20.0}{\celsius}}
\begin{tabular}{@{}cS[table-format=4.2]@{}}
\toprule
{$i$} & {$l_i$ (\si{\milli\meter})} \\
\midrule
1 & 3.92 \\
2 & 12.76 \\
3 & 21.60 \\
4 & 30.44 \\
5 & 39.28 \\
6 & 48.12 \\
7 & 56.96 \\
8 & 65.80 \\
9 & 74.64 \\
10 & 83.48 \\
\bottomrule
\end{tabular}
\end{table}

Frequency: $f = 39.875$ \si{\kilo\hertz}

Wavelength: $\lambda = 8.956$ \si{\milli\meter}

Speed: $v = 357.1$ \si{\meter\per\second}

\subsection{Sample Raw Data: P = 50\% (Atmospheric)}

Frequency: $f = 40.142$ \si{\kilo\hertz}

Wavelength: $\lambda = 8.934$ \si{\milli\meter}

Speed: $v = 358.6$ \si{\meter\per\second}

\subsection{Complete Pressure Data Summary}

(Full data available in pressure\_data.csv)

Key results:
\begin{itemize}
    \item Pressure range: 100\%--0.1\% of atmospheric
    \item Speed range: 356.8--370.7 \si{\meter\per\second}
    \item Average percentage error: 1.23\%
    \item Speed remains nearly constant with pressure (as expected for ideal gas)
    \item Small variations within experimental uncertainty
\end{itemize}

\section{Observations and Notes}

\subsection{Experimental Challenges}

\begin{itemize}
    \item \textbf{Temperature stability:} Maintaining constant temperature during measurements required careful control of hot plates. Some drift observed during long measurement sessions.
    \item \textbf{Vacuum chamber limitations:} Achieving very low pressures (< 1\%) was difficult and took considerable time. Leaks in chamber became noticeable at low pressures.
    \item \textbf{Measurement time:} Each temperature/pressure condition required $\sim$10--15 minutes for setup and 10 position measurements. Total experiment time: $\sim$8 hours.
    \item \textbf{Equipment calibration:} Verified frequency generator accuracy and vernier caliper zero before starting measurements.
\end{itemize}

\subsection{Data Quality Notes}

\begin{itemize}
    \item Linear fits consistently show $R^2 > 0.999$, indicating high-quality data
    \item Position measurements reproducible within \SI{0.02}{\milli\meter}
    \item Frequency remained stable during individual measurement sets
    \item Some temperature drift ($\sim 0.2$°C) observed during extended measurement sessions
\end{itemize}

\subsection{Additional Observations}

\begin{itemize}
    \item Speed clearly increases with temperature, following expected $\sqrt{T}$ relationship
    \item Pressure has minimal effect on speed (within experimental error), confirming ideal gas behavior
    \item Humidity correction important for accurate theoretical predictions
    \item Experimental values consistently within 1--2\% of theoretical predictions
\end{itemize}

\section{Sample Calculations}

\subsection{Example: Speed Calculation at T = 30.0°C}

Given:
\begin{itemize}
    \item Temperature: $t = 30.0$°C
    \item Frequency: $f = 40.253$ \si{\kilo\hertz}
    \item Wavelength (from fit): $\lambda = 9.045$ \si{\milli\meter}
\end{itemize}

Calculation:
\begin{align}
v &= f \lambda \\
&= (40.253 \times 10^3~\si{\hertz}) \times (9.045 \times 10^{-3}~\si{\meter}) \\
&= 363.9~\si{\meter\per\second}
\end{align}

Theoretical calculation:
\begin{align}
p_s &= 10^{10.286 - 1780/(237.3 + 30.0)} = 4243~\si{\pascal} \\
v_{\text{theory}} &= 331.5 \sqrt{(1 + 30.0/273.15)(1 + 0.16 \times 34 \times 4243/102400)} \\
&= 331.5 \sqrt{1.110 \times 1.0022} \\
&= 368.5~\si{\meter\per\second}
\end{align}

Percentage difference: $|363.9 - 368.5|/368.5 \times 100\% = 1.25\%$

\section{Data Files}

All raw data saved in CSV format:
\begin{itemize}
    \item \texttt{temperature\_data.csv} -- Summary of all temperature measurements
    \item \texttt{raw\_measurements\_temp.csv} -- Complete raw position data for temperature series
    \item \texttt{pressure\_data.csv} -- Summary of all pressure measurements
    \item \texttt{raw\_measurements\_pressure.csv} -- Complete raw position data for pressure series
\end{itemize}

\section{Acknowledgments}

We thank Prof. Xiawa Wang's lab for providing access to the vacuum chamber and environmental control equipment, which made this experiment possible.

\end{document}

